\subsection{HD-D03: small integral projectors and sharper constants}
\label{sec:small-projectors}

For integral maximal-order actions, the following explicit constants
replace the resultant constants in HD-D01:
\[
\begin{array}{c|r|r|r|r|r|r}
(n,K)&R&m&s&L&A_0&A_1\\ \hline
(2,\mathbb Q(\omega))&54&36&3&2&
2187&4455\\
(3,\mathbb Q(i))&768&400&64&8&
1048576&4634096\\
(3,\mathbb Q(\omega))&162&81&97200&486&
70858800&393697003\\
(4,\mathbb Q(\omega))&2624400&1920996&3&2&
35872267500&962620032542
\end{array}
\]
In each row every effective detector of degree \(D\) gives
\[
 c_t\le A_1\,\delta M_T D^2.
\]
These are exact integer constants. For an action of the order of
conductor \(f\), the same conclusion holds after multiplying \(A_1\)
by \(f^{8n}\). All cycle operations used below are integral
endomorphism pullbacks.

\begin{proof}
The proof of HD-D01 applies to any integral correspondence \(P\)
with \(P=R e_W\), whose positive and negative parts each have degree
at most \(mD\), together with an integral correspondence \(Z\)
satisfying
\[
 Z^\dagger=-Z,\qquad Z^2=-s\quad\hbox{on }W,
\]
whose weighted absolute coefficient sum is at most \(L\).
Indeed, with \(y=P[Y]\) and \(z=Zy\), the exact Gram matrix is
\(\operatorname{diag}(d,sd)\), the signed parts of \(z\) have
degrees at most \(LmD\), and every displayed line of the quadratic
extraction proof applies. Thus the two coefficients are
\[
 \left\lceil\frac{sR^2}{(n!)^2}\right\rceil,\qquad
 \left\lceil\frac{mR\sqrt{s(s+L^2)}}{n!}\right\rceil.           \tag{D03.1}
\]
The table evaluates these integers and has \(A_1\ge A_0\).
We now construct the actual correspondences.

\paragraph{Eisenstein fourfold.}
Write \(\Omega=[\omega]^*\), \(B=[1-\omega]^*\), and put
\[
 P=9(2-\Omega-\Omega^2)-(1-2\Omega+\Omega^2)B.
\]
On \(H^4\), the same cohomological calculation detailed for the
eightfold below gives
\(P=(2-\Omega-\Omega^2)(9-B-\bar B)\).
The first factor kills \(j=2\) and has eigenvalue three on
\(j=0,1,3,4\). The eigenvalues of \(B+\bar B\) are \(-9\)
at \(j=0,4\) and \(9\) at \(j=1,3\). Hence \(P=54e_W\).
Its two displayed parts have positive and negative degree sums
18 each, giving \(m=36\). Take \(Z=\Omega-\Omega^2\), with
\(s=3\) and \(L=2\). Formula (D03.1) gives
\(A_0=2187\), \(A_1=\lceil972\sqrt{21}\rceil=4455\).
Together with the principal Eisenstein detector bound \(D=300\),
this gives \(c_t\le400950000\,\delta M_T\).

\paragraph{Gaussian sixfold.}
Write \(I=[i]^*\), \(A=[2+i]^*\), and \(\bar A=[2-i]^*\), and set
\[
 P=(1-I)(150-A-\bar A).
\]
On the summand of \(H^6\) containing \(j\) factors from one complex
embedding, \(I\) has eigenvalue \(-(-1)^j\). It kills all odd \(j\).
At \(j=0,6\), \(A+\bar A\) has eigenvalue
\((2+i)^6+(2-i)^6=-234\); at \(j=2,4\), its eigenvalue is
\(25((2+i)^2+(2-i)^2)=150\). Consequently
\(P=2(150+234)e_W=768e_W\).
The explicit expansion
\[
 P=150+IA+I\bar A-150I-A-\bar A
\]
has positive and negative degree sums
\(150+125+125=400\).
Take \(Z=[1+i]^*\). Its endpoint eigenvalues on \(W\) are
\(-8i,8i\), giving \(Z^\dagger=-Z\) and \(Z^2=-64\).
Its degree multiplier on three-cycles is \(2^3=8\).

\paragraph{Eisenstein sixfold.}
Write \(\Omega=[\omega]^*\) and \(B=[1-\omega]^*\), and set
\[
 P=(1+\Omega+\Omega^2)(27-B).
\]
The unit sum selects precisely \(j=0,3,6\).
Since \((1-\omega)^6=-27\), the \(B\)-eigenvalues at these
three indices are respectively \(-27,27,-27\).
Therefore \(P=3\cdot54e_W=162e_W\).
The three positive terms \(27[\omega^k]^*\) and three negative
terms \([\,\omega^k(1-\omega)\,]^*\) each have total degree
\(3\cdot27=81\).
For \(A=[3+\omega]^*\), one has
\[
 (3+\omega)^6=37+360\omega=-143+180\sqrt{-3}.
\]
Thus \(Z=A+143\) is skew adjoint on \(W\), with
\(Z^2=-97200\), and its degree cost is \(7^3+143=486\).

\paragraph{Eisenstein eightfold.}
Retain \(\Omega\) and \(B\), and let \(A=[3+\omega]^*\),
\(\bar A=[3+\bar\omega]^*\). Use the explicit integral operator
\[
 E=2-\Omega-\Omega^2,\qquad
 F=81E-(1-2\Omega+\Omega^2)B,\qquad
 P=F(1127-A-\bar A).                                      \tag{D03.2}
\]
For verification of its spectrum, put \(\bar B=[1-\bar\omega]^*\).
The equality \(1-\bar\omega=-\omega^2(1-\omega)\), together with
\([-1]^*=1\) on \(H^8\), gives \(\bar B=\Omega^2B\) on this
cohomology group. The latter identity follows from
\([-1]^*=-1\) on \(H^1\) and exterior power functoriality.
It follows, on \(H^8\), that
\[
 F=E(81-B-\bar B).
\]
This is solely a cohomological verification of the explicitly
integral cycle operator in (D03.2); no equality on Chow groups is
needed.

The operator \(E\) kills \(j=1,4,7\) and has eigenvalue three
on the remaining indices. On \(j=0,2,6,8\), the eigenvalue of
\(B+\bar B\) is \(-81\); on \(j=3,5\), it is \(81\).
Thus \(F\) has eigenvalue 486 on \(j=0,2,6,8\) and zero
elsewhere. Finally \(A+\bar A\) has eigenvalue \(-4273\)
at \(j=0,8\), and eigenvalue \(1127\) at \(j=2,6\).
Hence
\[
 P=486(1127+4273)e_W=2624400e_W.
\]

The reduced expansion in (D03.2) is essential to the stated
degree bound. Its part \(81(2-\Omega-\Omega^2)\) has positive
and negative degree sums 162, and its part
\(-(1-2\Omega+\Omega^2)B\) has positive and negative degree
sums \(2\cdot3^4=162\). Thus \(F\) has each signed part bounded
by \(324D\). Multiplication by \(1127-A-\bar A\) gives
\[
 m=324(1127+2\cdot7^4)=1920996.
\]
Take \(Z=\Omega-\Omega^2\). On the two endpoint lines its
eigenvalues are \(\omega-\omega^2\) and its negative, so
\(Z^\dagger=-Z\), \(Z^2=-3\), and \(L=2\).

\paragraph{Exact ceilings and integral orders.}
For every row, the second integer in (D03.1) is the least
integer \(k\ge0\) satisfying
\[
 (k\,n!)^2\ge(mR)^2s(s+L^2).
\]
The companion exact script records both this inequality and the
strict reversed inequality for \(k-1\). There is no floating-point
rounding.

Finally, let the integral order be
\(\mathcal O_f=\mathbb Z+f\mathcal O_K\).
For every term \(c_u[u]^*\) in the displayed finite expansion
of \(P\), replace it by \(c_u[fu]^*\); do the same for \(Z\).
Every \(fu\) belongs to \(\mathcal O_f\), so these replacements
are actual integral endomorphism pullbacks on the given abelian
variety. On \(H^{2n}\),
\([fu]^*=f^{2n}[u]^*\), by the exterior power of the rational
\(H^1\) action. Therefore the new data satisfy
\[
 R_f=f^{2n}R,\quad m_f=f^{2n}m,\quad
 s_f=f^{4n}s,\quad L_f=f^{2n}L.
\]
Both unrounded expressions in (D03.1) are multiplied by
\(f^{8n}\); their integer ceilings are at most \(f^{8n}\)
times the original ceilings. This proves the asserted conductor
bound using explicit integral cycle maps.
\end{proof}

